\documentclass[11pt]{amsart}
\usepackage[utf8]{inputenc}
\usepackage[T1]{fontenc}
\usepackage{amssymb,amsmath}
\usepackage[margin=2.8cm]{geometry}
\usepackage{hyperref}

\newcommand{\stamptime}{\the\numexpr(\time-30)/60\relax:\ifnum\numexpr\time-60*((\time-30)/60)\relax<10 0\fi\the\numexpr\time-60*((\time-30)/60)\relax}
\AddToHook{shipout/foreground}{\put(0.5\paperwidth,-\paperheight+0.55cm){\makebox[0pt]{\tiny\ttfamily unit/doc/gcddomain.tex \ --- \ compiled \today, \stamptime}}}

\newtheorem{theorem}{Theorem}[section]
\newtheorem{lemma}[theorem]{Lemma}
\newtheorem{corollary}[theorem]{Corollary}
\theoremstyle{definition}
\newtheorem{definition}[theorem]{Definition}
\newtheorem{construction}[theorem]{Construction}
\theoremstyle{remark}
\newtheorem{remark}[theorem]{Remark}

% a Unit name
\newcommand{\nm}[1]{\textup{\texttt{#1}}}
% a Unit file
\newcommand{\file}[1]{\textup{\texttt{#1.seg}}}
\newcommand{\cont}{\mathrm{c}}
\newcommand{\gcdd}{\operatorname{gcd}}

\title[Gauss's lemma, file by file]{Gauss's lemma and polynomial gcd domains:\\
an informal proof following the Unit files}
\author{Thierry Coquand}
\date{4 November 2026}

\begin{document}
\maketitle

\section*{Purpose and how to read this text}

This text is written for a test. The claim is that the proof terms
of the checker Unit have the same structure as a detailed informal
proof. The example is the theorem
\begin{quote}
\emph{if $A$ is a gcd domain, then so is $A[X]$},
\end{quote}
together with Gauss's lemma, on which it rests. The development is in
the directory \url{unit/sag/gcddomain}. It has 23 files (plus
\file{All}, which only lists them) and checks in about $0.1$\,s.

\subsection*{Four levels of detail}
We follow van Benthem Jutting~\cite{jutting}, who developed a proof in
four stages: a proof~A ``which would normally be given in conversation'',
very informal and with gaps; a proof~B as in ``a rather pedantic
textbook''; a proof~C with more detail, where bars in the margin indicate
the context of assumptions and the small lemmas used are noted
explicitly; and a proof~D, the formal text checked by machine, with a
glossary for the human reader. Here:
\begin{itemize}
\item proofs A and B are in Part~\ref{part:AB}. Proof~B is close to the
  textbook of Mines, Richman and Ruitenburg~\cite{mrr}, and uses the
  same method as the formal proof;
\item proof~C is the rest of this text (Parts~\ref{part:base}
  and~\ref{part:algebra});
\item proof~D is the Unit files.
\end{itemize}
The modules of Unit play the role of Jutting's margin bars: a module
\nm{M (x : X) \{ \dots \}} fixes the assumptions $x\in X$ for the steps it
contains. His explicitly noted lemmas are our named small lemmas (such as
\nm{div-eq-right} or \nm{plus-assoc}), and his glossary is the
correspondence below. One difference: Jutting refined~A into~D, while
proof~C was written from the formal files. The question of the test is
how mechanical the passage from~C to~D is.

Proof~C has one section per file, in the order in which the files are
checked. Every definition, lemma and step is tagged with the name of the
corresponding Unit declaration, written like this: \nm{div-trans}. The
reader can therefore put the text and the file side by side.
Part~\ref{part:algebra} (seven files, about 900 lines) contains the
algebra. It is written in full: \emph{every} declaration of these files
appears in the text. Part~\ref{part:base} describes the files on rings and
polynomials more briefly, one paragraph per notion.

\subsection*{The correspondence}
The following constructs of Unit have a fixed informal reading.
\begin{itemize}
\item A \emph{declaration} \nm{name (x : X) : T = t} is a definition, a
  lemma or a step of a proof; its parameters are its hypotheses.
\item A \emph{module} \nm{module M (x : X) \{ \dots \}} is a proof, or a
  part of a text, under fixed hypotheses: ``Let $x\in X$. \dots''. Its
  declarations are the named steps of that proof, and \nm{M[a].step} refers
  to a step of it for the particular value $a$.
\item A \emph{record} is a structure ``consisting of \dots''. A
  $\Sigma$-type $(q : R) \times \dots$ is ``the datum of $q$ such
  that\dots''. A component of a pair is written \nm{h.1}, \nm{h.2}.
\item \nm{split} on a natural number, with the recursive value written
  \nm{| r}, is a proof by induction; \nm{r} is the induction hypothesis.
  \nm{split} on a decision $x = 0$ or $x \neq 0$ is a proof by cases.
  Unit writes a case analysis as its own small declaration, named after
  what is decided: \nm{step-byP} is ``case on whether $\cont(p)=0$''.
\item \nm{calc R \{ a =< h > b =< k > c \}} is a chain of equalities
  $a = b = c$, each step justified by the term between \nm{=<} and \nm{>}.
\item \nm{efq} is ``this case is impossible''.
\item An abbreviation \nm{module D = Domain[A,\dots]} means ``we use the
  results of the section on \nm{Domain} for this $A$''; they are then
  cited as \nm{D.cancel}, etc. The text cites them by their own names.
\end{itemize}
Some steps of the formal proof would be silent in a usual text. The most
frequent is \nm{div-eq-right}: ``$d \mid b$ and $b = b'$, hence
$d \mid b'$''. The text below leaves most of these steps silent, as a
usual proof does. They are the main reason why a formal step is often
longer than the corresponding sentence.

\subsection*{Foundations}
The development uses the base theory of Unit (dependent types, identity
types, inductive types and one universe per level). It adds no new
axioms. It rests on the parameters of the \nm{Theory} segment of
\file{All}: univalence (\nm{univ}) and replacement (\nm{images},
\nm{surjections}). They are used only to construct $A[X]$ as a quotient
of lists of coefficients, and to derive function extensionality.

\subsection*{Relation to Mines--Richman--Ruitenburg}
The proof of Gauss's lemma is the one of Mines, Richman and
Ruitenburg~\cite{mrr}, Chapter~IV, Lemma~4.3. It
differs in four points; Section~\ref{sec:gauss} states them. The theorem
on $A[X]$ is MRR Theorem~4.7(i), but its proof is different: MRR use the
field of fractions $K$ and gcds in $K[X]$. Here the Euclidean algorithm
in $K[X]$ is replaced by an elimination ``up to a non-zero scalar''
(Section~\ref{sec:elim}), and $K$ is not used.

\tableofcontents

%======================================================================
\part{Proofs A and B}\label{part:AB}

\section{Proof A}

\emph{Gauss's lemma: the content of $fg$ is $\cont(f)\cont(g)$.}
Dividing by the contents, it is enough to show that a product of two
primitive polynomials is primitive. Let $d$ divide all coefficients of
$fg$, and write $f = a + Xf_1$, $g = b + Xg_1$. Then $\gcdd(d,a)$ divides
all coefficients of $f_1g$, so by induction it divides $\cont(f_1)$,
hence all coefficients of $f$, so it is a unit. In the same way
$\gcdd(d,b)$ is a unit. Since $d\mid ab$, Euclid's lemma gives $d\mid b$,
so $d$ is a unit.

\emph{$A[X]$ is a gcd domain.} Run Euclid's algorithm on $f$ and $g$, but
multiply by non-zero scalars instead of dividing. This gives $D$ and
$\lambda\neq0$ with $D\mid\lambda f$, $D\mid\lambda g$, and every common
divisor of $f$ and $g$ divides $D$. By Gauss's lemma the primitive part
$E$ of $D$ divides $f$ and $g$, and $\gcdd(\cont(f),\cont(g))\,E$ is a
gcd of $f$ and $g$.

\section{Proof B}

Let $A$ be a gcd domain: an integral domain with $1\neq0$, with a
decision of $x=0$ for every $x$, and a chosen gcd $\gcdd(a,b)$ of every
pair. Gcds are determined up to a unit only, and divisibility $a\mid b$
means that $b = aq$ for some $q$.

\begin{lemma}\label{lemB:M}
(i) If $d\mid ua$ and $d\mid ub$ then $d\mid u\gcdd(a,b)$.
(ii) (Euclid) If $\gcdd(a,b)\mid1$ and $a\mid bc$ then $a\mid c$.
\end{lemma}
\begin{proof}
(i) If $u=0$ this is clear. Otherwise let $g = \gcdd(ua,ub)$. Then
$d\mid g$, and $u\mid g$, say $g = ut$. Cancelling $u$ in $ut\mid ua$
and $ut\mid ub$ gives $t\mid a$ and $t\mid b$, so $t\mid\gcdd(a,b)$ and
$d\mid g = ut\mid u\gcdd(a,b)$.
(ii) $a\mid ca$ and $a\mid cb$, so $a\mid c\gcdd(a,b)\mid c$ by~(i).
\end{proof}

The \emph{content} $\cont(f)$ of a polynomial $f$ is the gcd of its
coefficients, computed as
$\gcdd(f_n,\gcdd(f_{n-1},\dots\gcdd(f_0,0)))$. It divides every
coefficient of $f$, and every common divisor of the coefficients of $f$
divides it. If $\cont(f) = 0$ then $f = 0$. The polynomial $f$ is
\emph{primitive} if every common divisor of its coefficients divides~$1$.
If $\cont(f)\neq0$, then $f = \cont(f)f^*$ with $f^*$ primitive. By
Lemma~\ref{lemB:M}(i) and induction on the number of coefficients: if
$f$ is primitive and $d\mid uf_i$ for every $i$, then $d\mid u$.

\begin{theorem}\label{thmB:C}
If $d$ divides every coefficient of $fg$, then $d\mid\cont(f)\cont(g)$.
\end{theorem}
\begin{proof}
By induction on the total number $N$ of coefficients of $f$ and $g$, for
all $f$ and $g$. If $\cont(f) = 0$ or $\cont(g) = 0$, then
$\cont(f)\cont(g) = 0$ and there is nothing to prove; this covers $N=0$.
Otherwise write $f = \cont(f)f^*$ and $g = \cont(g)g^*$, with $f^*$ and
$g^*$ primitive and with as many coefficients as $f$ and $g$. Then
$fg = \cont(f)\cont(g)\,f^*g^*$, so it is enough to show that $f^*g^*$
is primitive.

So let $f = a + Xp$ and $g = b + Xq$ be primitive, and let $d$ divide
every coefficient of $fg$. Put $t = \gcdd(d,a)$. The coefficient of
$X^{i+1}$ in $fg$ is $aq_{i+1}+(pg)_i$. Since $t\mid a$, $t$ divides
every coefficient of $pg$. By the induction hypothesis, applied to $p$
and $g$ (we need the general statement here, since $p$ need not be
primitive), $t\mid\cont(p)\cont(g)$, and $\cont(g)\mid1$, so
$t\mid\cont(p)$. So $t$ divides $a$ and every coefficient of $p$, that
is, every coefficient of $f$. Hence $t\mid1$. In the same way
$\gcdd(d,b)\mid1$. Now $d\mid ab$, the constant coefficient of $fg$, so
$d\mid b$ by Lemma~\ref{lemB:M}(ii), and $d\mid\gcdd(d,b)\mid1$.
\end{proof}

\begin{corollary}[Gauss's lemma]\label{corB:gauss}
If $f$ is primitive and $d$ divides every coefficient of $fg$, then $d$
divides every coefficient of $g$. Hence, if $E$ is primitive,
$\lambda\neq0$ and $\lambda f = Eq$, then $E\mid f$.
\end{corollary}
\begin{proof}
$d\mid\cont(f)\cont(g)\mid\cont(g)$. For the second statement, $\lambda$
divides every coefficient of $\lambda f = Eq$, hence of $q$, so
$q = \lambda q'$; then $\lambda f = \lambda Eq'$, and $f = Eq'$.
\end{proof}

\begin{theorem}[elimination]\label{thmB:elim}
For $f,g\in A[X]$ there are $D\in A[X]$ and $\lambda\neq0$ in $A$ with
$D\mid\lambda f$, $D\mid\lambda g$, and such that every common divisor of
$f$ and $g$ divides $D$.
\end{theorem}
\begin{proof}
Let $f$ have degree $<n_f$ and $g$ degree $<n_g$. By induction on
$n_f+n_g$. If $n_f = 0$, then $f = 0$, and $D = g$, $\lambda = 1$; the
case $n_g = 0$ is symmetric. Otherwise let $\alpha$ and $\beta$ be the
coefficients of $f$ and $g$ of degree $n_f-1$ and $n_g-1$. If
$\alpha = 0$, $f$ has degree $<n_f-1$, and we apply the induction
hypothesis; likewise if $\beta = 0$. Otherwise, say $n_f\ge n_g$, and let
$k = n_f - n_g$. Then $f' = \beta f - \alpha X^kg$ has degree $<n_f-1$.
The induction hypothesis gives $(D,\lambda')$ for $f'$ and $g$. Take
$\lambda = \lambda'\beta\neq0$. Then
$\lambda f = \lambda'f' + \alpha X^k(\lambda'g)$ and
$\lambda g = \beta(\lambda'g)$ are divisible by $D$, and a common divisor
of $f$ and $g$ divides $f'$, hence $D$. If $n_f<n_g$, exchange $f$ and
$g$.
\end{proof}

\begin{theorem}
If $A$ is a gcd domain, then so is $A[X]$.
\end{theorem}
\begin{proof}
$f = 0$ if and only if $\cont(f) = 0$, which is decidable. If $fg = 0$,
then $\cont(f)\cont(g) = 0$ by Theorem~\ref{thmB:C}, so $f=0$ or $g=0$.
And $1\neq0$ in $A[X]$.

Gcds: if $f=0$ then $g$ is a gcd, and if $g = 0$ then $f$ is a gcd. Let
now $f,g\neq0$, and take $D$, $\lambda$ from Theorem~\ref{thmB:elim}.
Then $D\neq0$, since $D = 0$ would give $\lambda f = 0$. Write
$D = \cont(D)E$ with $E$ primitive. If $\lambda f = Dr$, then
$\lambda f = E(\cont(D)r)$, so $E\mid f$ by
Corollary~\ref{corB:gauss}; likewise $E\mid g$. Let
$\gamma = \gcdd(\cont(f),\cont(g))$ and $G = \gamma E$.

$G\mid f$: write $f = Eq$. Then $\gamma$ divides $\cont(f)$, hence every
coefficient of $Eq$, hence every coefficient of $q$ by
Corollary~\ref{corB:gauss}; so $q = \gamma s$ and $f = Gs$. Likewise
$G\mid g$.

If $h\mid f$ and $h\mid g$, then $h\mid D$, say $D = ht$, and $h\neq0$.
Write $h = \cont(h)h^*$ with $h^*$ primitive. Then $\cont(h)$ divides
$\cont(f)$ and $\cont(g)$, hence $\gamma$. And
$\cont(D)E = D = h^*(\cont(h)t)$, so $h^*\mid E$ by
Corollary~\ref{corB:gauss}. Hence $h = \cont(h)h^*\mid\gamma E = G$.
\end{proof}

\noindent Proof~B is a usual textbook proof. It treats elements of
$A[X]$ as polynomials and does not say how a gcd is chosen. The section
\emph{What a usual text leaves implicit}, at the end, lists what proof~C
adds to it.

%======================================================================
\part{Rings and polynomials}\label{part:base}
%======================================================================

\section{Logic, identity, functions (\file{Base}, \file{Prelude},
\file{Axioms}, \file{FunExtUA})}

\file{Base} declares the small data types: the empty type, the unit type
\nm{Top} with element \nm{tt}, sums \nm{Plus} with injections \nm{inl},
\nm{inr}, natural numbers, lists, and Boolean operations.

\file{Prelude} contains the elementary properties of equality, each by one
use of the elimination rule $J$: symmetry \nm{sym}, transitivity
\nm{trans}, congruence \nm{ap}, transport \nm{subst}, and congruence of a
binary operation \nm{cong2}: $a=a'$ and $b=b'$ give
$a\cdot b = a'\cdot b'$. A type is a \emph{proposition} (\nm{isProp}) if
any two of its elements are equal. It is a \emph{set} (\nm{isSet}) if all
its equality types are propositions. Two pairs are equal when their
first components are equal and their second components are equal up to
transport (\nm{Sigma-eq}). A $\Sigma$-type of propositions over a
proposition is a proposition (\nm{isProp-Sigma}).

\file{Axioms} states function extensionality and univalence as types
(\nm{FunExt}, \nm{Univalence}). It derives their consequences:
\nm{funext}, a product of propositions is a proposition
(\nm{isProp-Pi}), and a product of sets is a set (\nm{isSet-Pi}).
\file{FunExtUA} derives function extensionality from univalence (HoTT
book, \S4.9).

\section{Quotients (\file{Replacement}, \file{Truncation},
\file{Quotient})}

\file{Replacement} states the replacement principle of Cavallo and
Coquand: the image of a map from a small type into a locally small type is
small. \file{Truncation} defines the propositional truncation as the image
of $A\to\nm{Top}$. \file{Quotient} constructs, for a proposition-valued
equivalence relation $\sim$ on a type $P$, the quotient $P/{\sim}$
(\nm{Quot}). It consists of
\begin{itemize}
\item the class map \nm{inc}, with \nm{relates}: $p\sim q$ gives
  $\nm{inc}\,p = \nm{inc}\,q$, and \nm{set}: the quotient is a set;
\item \nm{induction}: \emph{to prove a proposition for every element of
  the quotient, it is enough to prove it for the classes $\nm{inc}\,p$};
\item \nm{Descend}: a map $P\to S$ into a set $S$ that respects $\sim$
  defines a map $P/{\sim}\to S$ (\nm{map}). Its value on $\nm{inc}\,p$ is
  the given one (\nm{beta}).
\end{itemize}
The principle \nm{induction} is used several times below. A
mathematician applies it without saying so when reasoning ``on a
representative''.

\section{Commutative rings and natural numbers (\file{Ring},
\file{NatOrder})}

\file{Ring} defines a commutative ring (\nm{CRing}) as a structure
consisting of a type \nm{R}, elements \nm{z} ($0$) and \nm{o} ($1$),
operations \nm{add}, \nm{mul}, \nm{neg}, and the laws: associativity
(\nm{addA}, \nm{mulA}), commutativity (\nm{addC}, \nm{mulC}), units
(\nm{addZ}, \nm{mulO}), inverse (\nm{addN}) and distributivity (\nm{dist}).
The module \nm{RL} proves the usual consequences, such as $0x=0$
(\nm{zero-mul}), $x0=0$ (\nm{mul-zero}), $1x=x$ (\nm{one-left}),
$u(-v)=-(uv)$ (\nm{mul-neg}), $-0=0$ (\nm{neg-zero}),
$(x+y)-y=x$ (\nm{add-neg-cancel}), $a(bc)=b(ac)$ (\nm{mul-swap}) and
$(ab)(cd)=(ac)(bd)$ (\nm{mul-four}). It also proves the identities on
natural numbers used for indices: \nm{plus-zero}, \nm{plus-suc},
\nm{plus-comm}, \nm{suc-inj} and \nm{suc-zero} ($n+1\neq0$).

\file{NatOrder} defines $m\le n$ (\nm{Le}) and $m<n$ (\nm{Lt}) by
recursion. It proves reflexivity, transitivity, the comparison
\nm{compare} ($m=n+d$ or $n=m+d$ for some $d$), and the case distinctions
\nm{le-or-gt} and \nm{le-cases}.

\section{Polynomials as lists (\file{Poly}, \file{Mul}, \file{RawLaws},
\file{RawMultiplication})}

A polynomial is represented by the list of its coefficients, the constant
term first (\nm{P}). Its \emph{length} is \nm{plen}, and its $k$-th
coefficient \nm{coef} is $0$ beyond the length. Two lists are
\emph{equivalent} (\nm{peq}) if they have the same coefficients, so
trailing zeros do not matter. Sum \nm{addP}, scalar multiple \nm{scal} and
product \nm{pmul} are defined by recursion on lists:
\[
  f\cdot 0 = 0,\qquad f\cdot(b + Xg) = bf + X(fg).
\]
Their coefficients are computed in \nm{coef-addP}, \nm{coef-scal} and
\nm{coef0-pmul} (the constant coefficient of $fg$ is $f_0g_0$). Beyond the
length every coefficient is zero (\nm{coef-beyond}).

\file{RawLaws} and \file{RawMultiplication} prove the ring laws up to
\nm{peq}, and that the operations respect \nm{peq}. Among them:
\nm{negP} and \nm{neg-coef} (negation), \nm{scal-cong}, \nm{scal-scal}
($a(bf) = (ab)f$), \nm{mul-cons-left} ($(a+Xp)q = aq + X(pq)$),
\nm{mul-cong}, \nm{mul-comm}, \nm{mul-scal-left} and \nm{mul-scal-right}.

\section{The ring $A[X]$ (\file{RingHom}, \file{PolynomialRing},
\file{Coefficients})}

\file{RingHom} defines ring homomorphisms; it is needed by
\file{PolynomialRing} for the constant map $A\to A[X]$.

\file{PolynomialRing} (module \nm{Extension}, for a ring $A$ whose
underlying type is a set) defines $A[X]$ as the quotient of lists by
\nm{peq} (\nm{Carrier}, with class map \nm{include}). Each operation
descends from lists to the quotient because it respects \nm{peq}
(\nm{addQ}, \nm{mulQ}, \nm{negQ}, \nm{zeroQ}, \nm{oneQ}). The ring laws
of the quotient follow from those on lists by \nm{forall2} and
\nm{forall3}: \emph{a proposition about two or three elements of $A[X]$
holds if it holds for lists}. The structure of commutative ring is
\nm{ring}. It also defines the constants \nm{constant}, the variable
\nm{variable} ($X$), $\nm{constant}(ab) = \nm{constant}(a)\nm{constant}(b)$
(\nm{constant-mul}), and the lemmas \nm{mul-comm}, \nm{mul-assoc},
\nm{mul-one}.

\file{Coefficients} (module \nm{Over}) defines the $k$-th coefficient of
an element of $A[X]$. It is obtained by \nm{Descend}, since equivalent
lists have the same coefficients (\nm{Coef}, \nm{coef}). On a class it is
the coefficient of the list (\nm{coef-include}). The file proves:
\begin{itemize}
\item \nm{coef-ext}: \emph{two polynomials with the same coefficients are
  equal};
\item \nm{law}: an equation between coefficients that holds on lists holds
  on $A[X]$, since equality of coefficients is a proposition;
\item the coefficients of a sum, a scalar multiple, a negative, a product
  by $X$, $0$ and a constant: \nm{coef-add}, \nm{coef-scale},
  \nm{coef-neg}, \nm{coef-variable} (with \nm{shifted} and
  \nm{shift-coef}), \nm{coef-zero}, \nm{coef-const0} and
  \nm{coef-constS};
\item $\nm{Bounded}\ d\ x$: every coefficient of $x$ from index $d$ on is
  zero, that is, $x$ has degree $<d$.
\end{itemize}

%======================================================================
\part{The algebra}\label{part:algebra}
%======================================================================

In this part \emph{every} declaration of the seven files is named.
Throughout, $A$ is a commutative ring with underlying type $R$.

\section{Divisibility and gcd domains (\file{GcdDomain})}

\begin{lemma}[\nm{empty-prop}, \nm{not-prop}]
The empty type is a proposition. For every type $X$, $\neg X$ is a
proposition.
\end{lemma}
\begin{proof}
An element of the empty type gives anything. $\neg X = X\to\emptyset$ is a
product of propositions (\nm{isProp-Pi}).
\end{proof}

\subsection{Divisibility (module \nm{Divisibility}, for a ring $A$)}

\begin{definition}[\nm{Div}]
$a\mid b$ is the datum of $q$ with $b = aq$.
\end{definition}

\noindent So a proof of $a\mid b$ is a pair (quotient, equation). We never
decide whether $a\mid b$.

\begin{lemma}\label{lem:div}
\begin{enumerate}
\item (\nm{div-refl}) $a\mid a$, with quotient $1$.
\item (\nm{div-zero}) $a\mid 0$, with quotient $0$.
\item (\nm{div-one}) $1\mid a$, with quotient $a$.
\item (\nm{div-trans}) If $a\mid b$ and $b\mid c$ then $a\mid c$.
\item (\nm{div-mul}, \nm{div-mul-left}) If $a\mid b$ then $a\mid bc$ and
  $a\mid cb$.
\item (\nm{div-factor}, \nm{div-factor-left}) $a\mid ac$ and $a\mid ca$,
  with quotient $c$.
\item (\nm{div-add}, \nm{div-neg}) If $a\mid b$ and $a\mid c$ then
  $a\mid b+c$; if $a\mid b$ then $a\mid -b$.
\item (\nm{div-scale}) If $a\mid b$ then $ca\mid cb$, with the same
  quotient.
\item (\nm{div-mul2}) If $a\mid b$ and $c\mid d$ then $ac\mid bd$.
\item (\nm{div-eq-left}, \nm{div-eq-right}) If $a\mid b$ and $a=a'$
  (resp.\ $b=b'$) then $a'\mid b$ (resp.\ $a\mid b'$).
\item (\nm{div-of-zero}) If $0\mid b$ then $b=0$.
\end{enumerate}
\end{lemma}
\begin{proof}
(4) If $b = ah$ and $c = bk$ then
$c = bk = (ah)k = a(hk)$, so the quotient is $hk$.
(5) If $b=ah$, the quotient is $hc$, since $bc = (ah)c = a(hc)$; for $cb$,
first $cb=bc$.
(6) $ac = ac$ by definition; $ca = ac$ by commutativity.
(7) If $b=ah$ and $c=ak$ then $b+c = ah+ak = a(h+k)$; and
$-b = -(ah) = a(-h)$.
(8) $cb = c(ah) = (ca)h$.
(9) If $b=ah$ and $d=ck$ then $bd = (ah)(ck) = (ac)(hk)$
(\nm{mul-four}).
(10) By transport along the equality.
(11) $b = 0\cdot q = 0$.
\end{proof}

\begin{definition}[\nm{GcdOf}]
A gcd of $a$ and $b$ is the datum of $d$ with $d\mid a$, $d\mid b$, and a
function giving $c\mid d$ for every $c$ with $c\mid a$ and $c\mid b$.
\end{definition}

\noindent A gcd is unique only up to a unit. So it is a datum, not a
property.

\subsection{Gcd domains}

\begin{definition}[\nm{GcdDomain}]
A \emph{(discrete) gcd domain} structure on $A$ consists of
\begin{itemize}
\item \nm{set}: $R$ is a set;
\item \nm{dec}: for every $x$, $x=0$ or $x\neq0$ (a decision);
\item \nm{integral}: if $x\neq0$ and $y\neq0$ then $xy\neq0$;
\item \nm{gcd}: a gcd of every pair $a,b$;
\item \nm{nontrivial}: $1\neq0$.
\end{itemize}
\end{definition}

\subsection{Discrete domains (module \nm{Domain}, for $A$ with
\nm{set}, \nm{dec}, \nm{integral})}

\begin{lemma}\label{lem:sub}
\begin{enumerate}
\item (\nm{mul-sub}) $u(a-b) = ua - ub$.
\item (\nm{sub-zero}) If $a - b = 0$ then $a = b$.
\item (\nm{sub-self}) $a - a = 0$.
\end{enumerate}
\end{lemma}
\begin{proof}
(1) By distributivity and $u(-b) = -(ub)$.
(2) $a = (a+b)-b = (b+a)-b = b+(a-b) = b + 0 = b$.
(3) This is the inverse law.
\end{proof}

\begin{lemma}[cancellation, \nm{cancel}]\label{lem:cancel}
If $u\neq0$ and $ua = ub$ then $a=b$.
\end{lemma}
\begin{proof}
We decide whether $a-b=0$ (\nm{cancel-cancelD}).
If $a-b = 0$, then $a=b$ by Lemma~\ref{lem:sub}(2).
If $a-b\neq0$, then $u(a-b)\neq0$ by integrality. But
$u(a-b) = ua - ub = ub - ub = 0$. This case is impossible.
\end{proof}

\begin{lemma}
\begin{enumerate}
\item (\nm{cancel-zero}) If $u\neq0$ and $ux=0$ then $x=0$.
\item (\nm{div-cancel}) If $u\neq0$ and $ua\mid ub$ then $a\mid b$, with
  the same quotient.
\item (\nm{div-prop}) If $a\neq0$ then $a\mid b$ is a proposition.
\item (\nm{nonzero-prop}) $x\neq0$ is a proposition.
\item (\nm{mul-nonzero}) If $x\neq0$ and $y\neq0$ then $xy\neq0$.
\end{enumerate}
\end{lemma}
\begin{proof}
(1) $ux = 0 = u0$; cancel $u$.
(2) If $ub = (ua)q$ then $ub = u(aq)$; cancel $u$.
(3) Two quotients $q$, $q'$ satisfy $aq = b = aq'$, so $q=q'$ by
cancellation. The two equations are then equal, since $R$ is a set.
(4) By \nm{not-prop}. (5) This is integrality.
\end{proof}

\subsection{Gcds and Lemma M (module \nm{Gcd}, for a gcd domain $A$)}

We write $\gcdd(a,b)$ for the chosen gcd (\nm{gcd}). Its three properties
are \nm{gcd-left} ($\gcdd(a,b)\mid a$), \nm{gcd-right}
($\gcdd(a,b)\mid b$) and \nm{gcd-univ} ($c\mid a$ and $c\mid b$ give
$c\mid\gcdd(a,b)$).

\begin{lemma}[Lemma M, \nm{lemma-M}]\label{lem:M}
If $d\mid ua$ and $d\mid ub$ then $d\mid u\gcdd(a,b)$.
\end{lemma}
\begin{proof}
We decide whether $u=0$ (\nm{lemma-M-byU}).

If $u = 0$ then $u\gcdd(a,b) = 0$, and $d\mid 0$.

If $u\neq0$ (module \nm{M}, with $u\neq0$, $d\mid ua$, $d\mid ub$ fixed):
\begin{itemize}
\item (\nm{g}) let $g = \gcdd(ua,ub)$;
\item (\nm{dg}) $d\mid g$, as a common divisor of $ua$ and $ub$;
\item (\nm{ug}) $u\mid g$, since $u\mid ua$ and $u\mid ub$;
\item (\nm{t}) let $t$ be the quotient, $g = ut$;
\item (\nm{ta}) $t\mid a$: since $ut = g\mid ua$, cancel $u$
  (\nm{div-cancel});
\item (\nm{tb}) $t\mid b$, in the same way;
\item (\nm{gu}) $g\mid u\gcdd(a,b)$: $t\mid\gcdd(a,b)$ by \nm{gcd-univ},
  hence $ut\mid u\gcdd(a,b)$ (\nm{div-scale}), and $ut = g$;
\item (\nm{result}) $d\mid g\mid u\gcdd(a,b)$.\qedhere
\end{itemize}
\end{proof}

\begin{lemma}[Euclid, \nm{euclid}]\label{lem:euclid}
If $\gcdd(a,b)\mid1$ and $a\mid bc$ then $a\mid c$.
\end{lemma}
\begin{proof}
$a\mid ca$ and $a\mid cb$ ($cb = bc$), so $a\mid c\gcdd(a,b)$ by
Lemma~\ref{lem:M} with $u=c$. And $c\gcdd(a,b)\mid c\cdot1 = c$
(\nm{div-scale}). So $a\mid c$.
\end{proof}

\subsection{Discrete fields (module \nm{FieldGcd})}

Let $F$ be a ring whose type is a set, in which every $x$ is $0$ or
invertible (\nm{dec}), and where $1\neq0$.

\begin{lemma}
\begin{enumerate}
\item (\nm{unit-nonzero}) An invertible element is not $0$.
\item (\nm{decZ}, by cases \nm{decZ-decD}) $x=0$ or $x\neq0$.
\item (\nm{integral}, by cases \nm{integral-intD}) If $x\neq0$ and
  $y\neq0$ then $xy\neq0$.
\item (\nm{unit-div}) An invertible element divides every element.
\item (\nm{gcd}, by cases \nm{gcd-gcdD}) Every pair has a gcd.
\end{enumerate}
Hence (\nm{domain}) $F$ is a gcd domain.
\end{lemma}
\begin{proof}
(1) If $xy=1$ and $x=0$ then $1 = 0\cdot y = 0$.
(2) If $x=0$, done; if $x$ is invertible, then $x\neq0$ by (1).
(3) Decide $x$. The case $x=0$ is impossible. If $xw=1$ and $xy=0$ then
$y = 0$ (\nm{unit-reg}).
(4) If $xy=1$ then $b = 1b = (xy)b = x(yb)$.
(5) Decide $a$. If $a=0$, then $b$ is a gcd: $b\mid0=a$, $b\mid b$, and
a common divisor divides $b$. If $a$ is invertible, then $a$ is a gcd:
$a\mid a$, $a\mid b$ by (4), and a common divisor divides $a$.
\end{proof}

\section{The content of a sequence (\file{Content})}

Let $A$ be a gcd domain (module \nm{Over}). We work with sequences
$e\colon\mathbb N\to R$.

\begin{definition}
\begin{itemize}
\item (\nm{Below}) A property holds \emph{below $n$} if it holds for all
  $i<n$.
\item (\nm{Bnd}) $e$ is \emph{bounded by $n$} if $e_i = 0$ for $i\ge n$.
  Formally: $e_{n+i}=0$ for all $i$.
\item (\nm{Prim}) $e$ is \emph{primitive} if every common divisor of the
  $e_i$ divides $1$.
\item (\nm{cont}) $\cont_0(e) = 0$ and
  $\cont_{n+1}(e) = \gcdd(e_n, \cont_n(e))$.
\end{itemize}
\end{definition}

\noindent Two facts on indices: $n<n+1$ (\nm{lt-top}), and $i<n$ gives
$i<n+1$ (\nm{lt-up}).

\begin{lemma}\label{lem:cont}
\begin{enumerate}
\item (\nm{cont-div}) $\cont_n(e)\mid e_i$ for $i<n$.
\item (\nm{cont-univ}) If $d\mid e_i$ for all $i<n$, then
  $d\mid\cont_n(e)$.
\item (\nm{cont-M}) If $d\mid ue_i$ for all $i<n$, then
  $d\mid u\cont_n(e)$.
\end{enumerate}
\end{lemma}
\begin{proof}
All three are proved by induction on $n$.

(1) For $n=0$ there is no $i<0$. For $n+1$, let $i<n+1$; then $i<n$ or
$i=n$ (\nm{cont-div-step}). If $i<n$:
$\cont_{n+1}(e) = \gcdd(e_n,\cont_n(e))\mid\cont_n(e)\mid e_i$, by the
induction hypothesis. If $i=n$: $\gcdd(e_n,\cont_n(e))\mid e_n$.

(2) For $n=0$: $d\mid0$. For $n+1$: $d\mid e_n$, and $d\mid\cont_n(e)$
by the induction hypothesis, so $d$ divides their gcd.

(3) For $n=0$: $d\mid0 = u\cdot0$. For $n+1$: $d\mid ue_n$, and
$d\mid u\cont_n(e)$ by the induction hypothesis. Lemma~\ref{lem:M} gives
$d\mid u\gcdd(e_n,\cont_n(e))$.
\end{proof}

\subsection{Bounded sequences (module \nm{Bounded}, for $e$ bounded by
$n$)}

\begin{lemma}\label{lem:bounded}
Let $e$ be bounded by $n$, and let $\cont(e) = \cont_n(e)$ (\nm{c}).
\begin{enumerate}
\item (\nm{every}) If $d\mid e_i$ for all $i<n$, then $d\mid e_i$ for
  all $i$.
\item (\nm{c-div}) $\cont(e)\mid e_i$ for all $i$.
\item (\nm{c-univ}) If $d\mid e_i$ for all $i$ then $d\mid\cont(e)$.
\item (\nm{c-zero}) If $\cont(e) = 0$ then $e_i = 0$ for all $i$.
\item (\nm{prim-c}) If $e$ is primitive then $\cont(e)\mid1$.
\item (Corollary, \nm{prim-M}) If $e$ is primitive and $d\mid ue_i$ for
  all $i$, then $d\mid u$.
\end{enumerate}
\end{lemma}
\begin{proof}
(1) Either $n\le i$ or $i<n$ (\nm{every-step}). If $n\le i$, then
$i = n + k$ for some $k$, so $e_i = 0$ and $d\mid0$. If $i<n$, use the
hypothesis.
(2) By (1) and Lemma~\ref{lem:cont}(1).
(3) By Lemma~\ref{lem:cont}(2).
(4) $0 = \cont(e)\mid e_i$, so $e_i=0$ (\nm{div-of-zero}).
(5) $\cont(e)$ is a common divisor of the $e_i$, by (2).
(6) $d\mid u\cont(e)$ by Lemma~\ref{lem:cont}(3), and
$u\cont(e)\mid u\cdot1 = u$ by (5).
\end{proof}

\section{Polynomials over a gcd domain (\file{PolyBase})}

Let $A$ be a gcd domain (module \nm{Over}). We write \nm{Pol} for $A[X]$
and \nm{coef} for the coefficients of its elements. For $a\in R$ and
$x\in A[X]$, $ax$ (\nm{sc}) is the product of the constant $a$ with $x$.

\begin{lemma}
\begin{enumerate}
\item (\nm{coef-sc}) $(ax)_k = a x_k$.
\item (\nm{coef-mul-include}) The coefficients of the product of two
  classes are those of the product of the lists.
\item (\nm{bnd-list}) The coefficients of a list are bounded by its
  length.
\end{enumerate}
\end{lemma}

\begin{definition}[\nm{Prim}]
$x\in A[X]$ is \emph{primitive} if its sequence of coefficients is
primitive. If the class of a list is primitive, so is the list
(\nm{prim-include}).
\end{definition}

\subsection{Division by a scalar}

\begin{construction}[\nm{divL}, \nm{divL-coef}]
Let $q$ be a list and suppose that $a$ divides each coefficient of $q$.
The list $q/a$ is obtained by replacing each coefficient by its quotient,
and $q_i = a(q/a)_i$ for every $i$.
\end{construction}
\begin{proof}
By induction on the list. The empty list gives the empty list, and
$q_i = 0 = a\cdot0$. For $c + Xq$, the head becomes the quotient of $c$,
and the tail is obtained by the induction hypothesis.
\end{proof}

\noindent Let now $a\neq0$ be fixed (module \nm{Scalar}). For $x\in A[X]$,
\nm{Quot}~$x$ is the datum of $y$ with $x = ay$, and \nm{Divides}~$x$ means
that $a$ divides every coefficient of $x$.

\begin{lemma}
\begin{enumerate}
\item (\nm{cancel}) If $ay = ay'$ then $y=y'$.
\item (\nm{quot-prop}) \nm{Quot}~$x$ is a proposition.
\item (\nm{divide}) If $a$ divides every coefficient of $x$, there is $y$
  with $x = ay$.
\end{enumerate}
\end{lemma}
\begin{proof}
(1) For every $k$, $ay_k = (ay)_k = (ay')_k = ay'_k$; cancel $a$
(Lemma~\ref{lem:cancel}). Then $y=y'$ by \nm{coef-ext}.

(2) Two quotients are equal by (1). The equations are then equal, since
$A[X]$ is a set.

(3) Since \nm{Quot}~$x$ is a proposition, it is enough to prove it when
$x$ is the class of a list $q$ (\nm{induction}). In that case
(\nm{divide-list}), $a$ divides each coefficient of the list $q$ (step
\nm{w'}). Take $y$ the class of $q/a$; then $x = ay$ coefficientwise, by
\nm{divL-coef}.
\end{proof}

\section{Gauss's lemma (\file{Gauss})}\label{sec:gauss}

The proof is MRR Lemma~IV.4.3. It differs in four points.
\begin{enumerate}
\item MRR remove the leading term $f_mX^m$. Here lists start with the
  constant term, so we remove the constant term: $f = a + Xf_1$.
\item MRR state $\cont(fg) = \cont(f)\cont(g)$. We prove the statement
  (C) below (one divisibility), and get the other one directly.
\item MRR say ``we may divide $f$ and $g$ by their contents and prove the
  claim for primitive $f$ and $g$''. But the induction hypothesis is then
  applied to $f - f_mX^m$, which is not primitive, so it must be the
  general statement. Here this is explicit. The induction is on the total
  length $N$, and the statement is the general one (C). At each length we
  first prove the primitive case from (C) at the length below
  (\nm{PrimStep}), then (C) by dividing by the contents (\nm{Nonzero}).
  Division by a scalar keeps the length, so no second use of the
  induction hypothesis is needed.
\item MRR conclude with their Lemma~1.1(iii). We use Euclid
  (Lemma~\ref{lem:euclid}), which is their Lemma~1.1(iv).
\end{enumerate}

\subsection{On lists (module \nm{Direct})}

Let $A$ be a gcd domain. In this subsection polynomials are lists.

\begin{definition}
\begin{itemize}
\item (\nm{content}) The content of a list $p$ is
  $\cont(p) = \cont_{n}(p)$, where $n$ is the length of $p$. Then
  Lemma~\ref{lem:bounded} applies to $p$ (module \nm{At}, with \nm{K}).
\item (\nm{Common}) $d$ is a \emph{common divisor} of $p$ if it divides
  every coefficient of $p$.
\item (\nm{Product}) Statement (C) for $p$, $q$: every common divisor
  of $pq$ divides $\cont(p)\cont(q)$.
\end{itemize}
\end{definition}

\begin{lemma}
\begin{enumerate}
\item (\nm{div-sub}) If $d\mid a$ and $d\mid b$ then $d\mid a-b$.
\item (\nm{sub-add}) $(x+y)-x = y$.
\end{enumerate}
\end{lemma}

\begin{lemma}[\nm{tail-common}]\label{lem:tail}
If $d\mid a$ and $d$ is a common divisor of $(a+Xp)q$, then $d$ is a
common divisor of $pq$.
\end{lemma}
\begin{proof}
Fix $i$ (module \nm{Tail}). Let $x = aq_{i+1}$ (\nm{x}) and
$y = (pq)_i$ (\nm{y}).
\begin{itemize}
\item (\nm{eq}) $((a+Xp)q)_{i+1} = x + y$, since $(a+Xp)q = aq + X(pq)$;
\item (\nm{dx}) $d\mid x$, since $d\mid a$;
\item (\nm{dxy}) $d\mid x+y$, by hypothesis and (\nm{eq});
\item (\nm{result}) $d\mid (x+y) - x = y$.\qedhere
\end{itemize}
\end{proof}

\begin{definition}[\nm{Rec}]
$\mathrm{Rec}(N)$: statement (C) holds for all $p,q$ of total length $N$.
\end{definition}

\begin{lemma}[primitive step, \nm{PrimStep}]\label{lem:primstep}
Assume $\mathrm{Rec}(N)$. Let $a+Xp$ and $b+Xq$ be primitive, of total
length $N+1$, and let $d$ be a common divisor of their product. Then
$d\mid1$.
\end{lemma}
\begin{proof}
Write $f = a+Xp$ and $g = b+Xq$. First, $p$ and $g$ have total length $N$
(\nm{ef}), and so do $q$ and $f$ (\nm{es}).

\emph{Claim} (module \nm{One}). Let $a_0+Xp_0$ be primitive and $q_0$
primitive, with $p_0$, $q_0$ of total length $N$, and let $d$ be a common
divisor of $(a_0+Xp_0)q_0$. Then $\gcdd(d,a_0)\mid1$.
\begin{itemize}
\item (\nm{t}) Let $t=\gcdd(d,a_0)$.
\item (\nm{ht}) $t$ is a common divisor of $(a_0+Xp_0)q_0$, since
  $t\mid d$.
\item (\nm{tc}) $t\mid\cont(p_0)\cont(q_0)$: by Lemma~\ref{lem:tail}
  ($t\mid a_0$), $t$ is a common divisor of $p_0q_0$; apply
  $\mathrm{Rec}(N)$.
\item (\nm{ct}) $t\mid\cont(p_0)$, since
  $\cont(p_0)\cont(q_0)\mid\cont(p_0)\cdot1$, because $q_0$ is primitive
  (\nm{prim-c}).
\item (\nm{coeffs}) $t$ is a common divisor of $a_0+Xp_0$: the
  coefficient of index $0$ is $a_0$, and $t\mid a_0$; the coefficient of
  index $i+1$ is $(p_0)_i$, and $t\mid\cont(p_0)\mid(p_0)_i$.
\item (\nm{result}) $t\mid1$, since $a_0+Xp_0$ is primitive.
\end{itemize}

Now:
\begin{itemize}
\item (\nm{cpA}) $\gcdd(d,a)\mid1$, by the claim for $a$, $p$, $g$.
\item (\nm{hs}) $d$ is a common divisor of $gf$, since $gf$ and $fg$ have
  the same coefficients.
\item (\nm{cpB}) $\gcdd(d,b)\mid1$, by the claim for $b$, $q$, $f$.
\item (\nm{ab}) $d\mid ab$, the constant coefficient of $fg$.
\item (\nm{db}) $d\mid b$, by Euclid (Lemma~\ref{lem:euclid}) and
  (\nm{cpA}).
\item (\nm{result}) $d\mid\gcdd(d,b)\mid1$, by (\nm{db}) and
  (\nm{cpB}).\qedhere
\end{itemize}
\end{proof}

\begin{lemma}[\nm{not-prim-nil}, \nm{primitive}]\label{lem:primitive}
The empty list is not primitive. Hence, under $\mathrm{Rec}(N)$, the
product of two primitive lists of total length $N+1$ is primitive.
\end{lemma}
\begin{proof}
$0$ is a common divisor of the empty list. If it were primitive, $0\mid1$,
so $1=0$, against \nm{nontrivial}. For the second statement, by cases on
the two lists: if one is empty, it is not primitive, so the case is
impossible. Otherwise they are $a+Xp$ and $b+Xq$, and we apply
Lemma~\ref{lem:primstep}.
\end{proof}

\begin{lemma}[primitive part of a list, \nm{Part}]\label{lem:part}
Let $p$ be a list with $\cont(p)\neq0$, and let $c = \cont(p)$ (\nm{c}).
Let $p^* = p/c$ (\nm{qs}). Then
\begin{enumerate}
\item (\nm{length}) $p^*$ has the same length as $p$ (\nm{div-len});
\item (\nm{eq}, \nm{scaled}) $p_i = cp^*_i$ for every $i$, so $p$ is
  equivalent to $cp^*$;
\item (\nm{prim}) $p^*$ is primitive.
\end{enumerate}
\end{lemma}
\begin{proof}
(1) Division replaces each coefficient and keeps the list structure (by
induction, \nm{div-len}). (2) By \nm{divL-coef}.
(3) Let $d$ be a common divisor of $p^*$. Then $cd\mid cp^*_i = p_i$ for
every $i$, so $cd\mid c = c\cdot1$ (\nm{c-univ}). Cancel $c\neq0$:
$d\mid1$.
\end{proof}

\begin{lemma}[non-zero contents, \nm{Nonzero}]\label{lem:nonzero}
Assume $\mathrm{Rec}(N)$. Let $p,q$ have total length $N+1$ and non-zero
contents. Then (C) holds for $p,q$.
\end{lemma}
\begin{proof}
Take the primitive parts $p^*$, $q^*$ (\nm{Pp}, \nm{Pq}), and let
$u = \cont(p)\cont(q)$ (\nm{u}).
\begin{itemize}
\item (\nm{len'}) $p^*$ and $q^*$ have total length $N+1$, by
  Lemma~\ref{lem:part}(1).
\item (\nm{prim}) $p^*q^*$ is primitive, by Lemma~\ref{lem:primitive}.
\item (\nm{scaled}) $pq$ is equivalent to $u(p^*q^*)$:
  \[pq \sim (\cont(p)p^*)(\cont(q)q^*) \sim \cont(p)(p^*(\cont(q)q^*))
  \sim \cont(p)(\cont(q)(p^*q^*)) \sim u(p^*q^*).\]
\item (\nm{result}) Let $d$ be a common divisor of $pq$. Then
  $d\mid (pq)_i = u(p^*q^*)_i$ for every $i$. Since $p^*q^*$ is primitive,
  $d\mid u$ (\nm{prim-M}).\qedhere
\end{itemize}
\end{proof}

\begin{theorem}[statement (C), \nm{product}]\label{thm:C}
For all lists $p,q$, every common divisor of $pq$ divides
$\cont(p)\cont(q)$.
\end{theorem}
\begin{proof}
By induction on the total length $N$ (\nm{induction}).

$N=0$ (\nm{at-zero}): $p$ is empty, so $\cont(p)\cont(q) = 0$ and $d\mid0$.
(If $p$ were not empty, its length would be positive.)

$N+1$ (\nm{step}), with the induction hypothesis $\mathrm{Rec}(N)$:
decide whether $\cont(p) = 0$ (\nm{step-byP}).
If $\cont(p)=0$ then $\cont(p)\cont(q)=0$, and $d\mid0$.
Otherwise, decide whether $\cont(q)=0$ (\nm{step-byQ}). If $\cont(q)=0$,
likewise. Otherwise apply Lemma~\ref{lem:nonzero}.
\end{proof}

\begin{lemma}[\nm{multiply-common}]\label{lem:multcommon}
If $c$ is a common divisor of $p$ and $d$ a common divisor of $q$, then
$cd$ is a common divisor of $pq$.
\end{lemma}
\begin{proof}
By induction on $q$. For the empty list, $pq$ is empty. For $b+Xq$,
$p(b+Xq) = bp + X(pq)$. The coefficient $i$ of $bp$ is $p_ib$, and
$cd\mid p_ib$ (\nm{div-mul2}). The coefficients of $X(pq)$ are $0$ and
those of $pq$, which $cd$ divides by the induction hypothesis
(\nm{multiply-common-shift-common}).
\end{proof}

\begin{corollary}[multiplicativity of the content]
$\cont(pq)\mid\cont(p)\cont(q)$ (\nm{content-mul-left}, by
Theorem~\ref{thm:C}), and $\cont(p)\cont(q)\mid\cont(pq)$
(\nm{content-mul-right}, by Lemma~\ref{lem:multcommon} and \nm{c-univ}).
\end{corollary}

\begin{corollary}[Gauss's lemma on lists, \nm{gauss}]\label{cor:gauss}
If $p$ is primitive and $d$ is a common divisor of $pq$, then $d$ is a
common divisor of $q$.
\end{corollary}
\begin{proof}
$d\mid\cont(p)\cont(q)$ by Theorem~\ref{thm:C}, and
$\cont(p)\cont(q)\mid\cont(q)\cdot 1$ since $\cont(p)\mid1$. So
$d\mid\cont(q)\mid q_i$.
\end{proof}

\subsection{On $A[X]$ (module \nm{Over})}

Corollary~\ref{cor:gauss} is restated as \nm{gauss-list}.

\begin{theorem}[Gauss's lemma, \nm{GaussAt}]\label{thm:gauss}
Let $d\neq0$. If $x\in A[X]$ is primitive and $d$ divides every
coefficient of $xy$, then $d$ divides every coefficient of $y$.
\end{theorem}
\begin{proof}
Let \nm{Concl}~$y$ be the conclusion, and \nm{Stmt}~$x\,y$ the whole
implication. Since $d\neq0$, each $d\mid y_j$ is a proposition
(\nm{div-prop}), so \nm{Concl}~$y$ (\nm{concl-prop}) and \nm{Stmt}~$x\,y$
(\nm{stmt-prop}) are propositions. By \nm{forall2} it is enough to prove
the statement for classes of lists (\nm{on-lists}). For these the
coefficients are those of the lists, and we apply \nm{gauss-list}.
\end{proof}

\begin{remark}
A usual text would not have the hypothesis $d\neq0$. It is needed only to
pass from lists to $A[X]$. For $d=0$, a proof of $0\mid0$ may have any
quotient, so $0\mid y_j$ is not a proposition, and the argument ``it is
enough to prove it for lists'' does not apply. Where $d=0$ is needed
(integrality, below), the statement on lists is used directly.
\end{remark}

\begin{lemma}[division by a primitive, \nm{DivPrim}]\label{lem:divprim}
Let $E$ be primitive, $\lambda\neq0$ and $\lambda f = Eq$. Then $E\mid f$.
\end{lemma}
\begin{proof}
(With \nm{S} the division by $\lambda$.)
\begin{itemize}
\item (\nm{dq}) $\lambda$ divides every coefficient of $q$: it divides
  every coefficient of $\lambda f = Eq$; apply Theorem~\ref{thm:gauss}.
\item (\nm{q'}) So $q = \lambda q'$ for some $q'$.
\item (\nm{same}) $\lambda f = Eq = E(\lambda q') = \lambda(Eq')$.
\item (\nm{result}) $f = Eq'$, by cancelling $\lambda$.\qedhere
\end{itemize}
\end{proof}

\section{Elimination (\file{Elimination})}\label{sec:elim}

This replaces the Euclidean algorithm in $K[X]$ used by MRR. We stay in
$A[X]$ and allow multiplication by a non-zero scalar.

\begin{lemma}[two ring identities, module \nm{Ident}]
In every commutative ring:
\begin{enumerate}
\item (\nm{ident}) $(LB)f = L(Bf - Tg) + T(Lg)$;
\item (\nm{ident2}) $(LB)g = B(Lg)$.
\end{enumerate}
\end{lemma}
\begin{proof}
(1) $L(Bf-Tg) + T(Lg) = (L(Bf) + L(-(Tg))) + T(Lg)
= (L(Bf) - T(Lg)) + T(Lg) = L(Bf) + (-(T(Lg)) + T(Lg)) = L(Bf) + 0
= L(Bf) = (LB)f$.
(2) $(LB)g = L(Bg) = B(Lg)$.
\end{proof}

\noindent Let now $A$ be a gcd domain (module \nm{Over}), and write
$\nm{Bounded}\ n\ x$ for ``$x$ has degree $<n$'' (\nm{Bounded}).

\begin{definition}[\nm{Elim}]
An \emph{elimination datum} for $f,g\in A[X]$ consists of $D\in A[X]$
(\nm{gd}) and $\lambda\in A$ (\nm{lam}) with $\lambda\neq0$
(\nm{lam-nz}), $D\mid\lambda f$ (\nm{div-f}), $D\mid\lambda g$
(\nm{div-g}), and such that every common divisor of $f$ and $g$ divides
$D$ (\nm{common}).
\end{definition}

\noindent An elimination datum for $f,g$ is one for $g,f$ (\nm{swap}).

\begin{lemma}[\nm{sc-one}, \nm{zero-left}]\label{lem:zeroleft}
$1x = x$. If $f$ has degree $<0$ then $(g,1)$ is an elimination datum for
$f,g$.
\end{lemma}
\begin{proof}
$f=0$, since all its coefficients are $0$ (\nm{f0}, by \nm{coef-ext}).
Then $g\mid0 = 1f$, $g\mid g = 1g$, and a common divisor of $f$, $g$
divides $g$.
\end{proof}

\begin{lemma}[indices and powers of $X$]
\begin{enumerate}
\item (\nm{plus-assoc}) $(a+b)+c = a+(b+c)$.
\item (\nm{xp}) $X^kx$ is defined by recursion on $k$.
\item (\nm{coef-xp}) $(X^kx)_{k+i} = x_i$.
\item (\nm{xp-mul}) $X^kx = (X^k\cdot1)x$.
\item (\nm{bound-down}) If $x$ has degree $<n+1$ and $x_n=0$, then $x$
  has degree $<n$.
\end{enumerate}
\end{lemma}
\begin{proof}
(1)--(4) By induction on $a$, resp.\ $k$. (5) The coefficient $n+0$ is
$x_n = 0$, and the coefficient $n+(i+1) = (n+1)+i$ is $0$ by hypothesis.
\end{proof}

\noindent The identities of \nm{Ident} are used in the ring $A[X]$
(\nm{IB}).

\begin{definition}[\nm{Rec}]
$\mathrm{Rec}(N)$: for all $f,g$ and $n_f,n_g$ with $n_f+n_g = N$, $f$ of
degree $<n_f$ and $g$ of degree $<n_g$, there is an elimination datum for
$f,g$.
\end{definition}

\begin{lemma}[the elimination step, \nm{Reduce}]\label{lem:reduce}
Assume $\mathrm{Rec}(N)$. Let $f$ have degree $<n'_f+1$ and $g$ degree
$<n'_g+1$, with $n'_f = k + n'_g$ and $n'_f + (n'_g+1) = N$. Assume that
the top coefficient $g_{n'_g}$ is non-zero. Then $f,g$ have an
elimination datum.
\end{lemma}
\begin{proof}
Let $\alpha = f_{n'_f}$ (\nm{alpha}), $\beta = g_{n'_g}$ (\nm{beta}),
$T = \alpha X^k$ (\nm{T}), and
\[
  f' = \beta f - Tg\qquad(\nm{f'}).
\]
\begin{itemize}
\item (\nm{coef-T}) $(Tg)_j = \alpha(X^kg)_j$.
\item (\nm{coef-f'}) $f'_j = \beta f_j - \alpha(X^kg)_j$.
\item (\nm{idx}) $(X^kg)_{n'_f+i} = g_{n'_g+i}$, since
  $n'_f + i = k + (n'_g+i)$.
\item (\nm{top}) $f'_{n'_f} = \beta\alpha - \alpha\beta = 0$.
\item (\nm{above}) $f'_{n'_f+j+1} = \beta\cdot0 - \alpha\cdot0 = 0$,
  by the bounds on $f$ and $g$.
\item (\nm{bounded}) So $f'$ has degree $<n'_f$, by (\nm{top}) and
  (\nm{above}).
\item (\nm{common'}) A common divisor of $f$ and $g$ divides
  $\beta f - Tg = f'$.
\item (\nm{rec}) Since $n'_f + (n'_g+1) = N$, $\mathrm{Rec}(N)$ gives an
  elimination datum $(D,\lambda')$ for $f',g$.
\item (\nm{lam}) Let $\lambda = \lambda'\beta$.
\item (\nm{sc-f}) $\lambda f = \lambda'f' + T(\lambda'g)$, by
  \nm{ident}.
\item (\nm{sc-g}) $\lambda g = \beta(\lambda'g)$, by \nm{ident2}.
\item (\nm{result}) $(D,\lambda)$ is an elimination datum for $f,g$:
  $\lambda\neq0$ since $\lambda'\neq0$ and $\beta\neq0$;
  $D\mid\lambda'f'$ and $D\mid\lambda'g$, so $D$ divides $\lambda f$ and
  $\lambda g$; and a common divisor of $f,g$ divides $f'$ and $g$, hence
  $D$.\qedhere
\end{itemize}
\end{proof}

\begin{lemma}[both bounds positive, \nm{Tops}]\label{lem:tops}
Assume $\mathrm{Rec}(N)$. Let $f$ have degree $<n'_f+1$ and $g$ degree
$<n'_g+1$, with $(n'_f+1)+(n'_g+1) = N+1$. Then $f,g$ have an elimination
datum.
\end{lemma}
\begin{proof}
Let $\alpha = f_{n'_f}$ and $\beta = g_{n'_g}$. The arithmetic:
$n'_f + (n'_g+1) = N$ (\nm{eF}), $(n'_f+1)+n'_g = N$ (\nm{eG}) and
$n'_g + (n'_f+1) = N$ (\nm{eSym}).

Decide whether $\alpha = 0$ (\nm{byF}). If $\alpha=0$, $f$ has degree
$<n'_f$ (\nm{bound-down}); apply $\mathrm{Rec}(N)$ with (\nm{eF}).

Otherwise decide whether $\beta = 0$ (\nm{byG}). If $\beta = 0$, $g$ has
degree $<n'_g$; apply $\mathrm{Rec}(N)$ with (\nm{eG}).

Otherwise compare $n'_f$ and $n'_g$ (\nm{byCmp}). If $n'_f = n'_g + d$,
apply Lemma~\ref{lem:reduce} with $k = d$. If $n'_g = n'_f + d$, apply
Lemma~\ref{lem:reduce} to $g, f$ (with (\nm{eSym}) and $\alpha\neq0$),
and exchange (\nm{swap}).
\end{proof}

\begin{theorem}[elimination, \nm{elimination}]\label{thm:elim}
If $f$ has degree $<n_f$ and $g$ has degree $<n_g$, then $f,g$ have an
elimination datum.
\end{theorem}
\begin{proof}
We prove $\mathrm{Rec}(N)$ by induction on $N$ (\nm{elim}), and apply it
with $N = n_f+n_g$.

$N = 0$ (\nm{onZero}): then $n_f = 0$, and Lemma~\ref{lem:zeroleft}
applies. ($n_f$ cannot be a successor, since $n_f+n_g = 0$.)

$N+1$: by cases on $n_f$ (\nm{onF}). If $n_f = 0$, apply
Lemma~\ref{lem:zeroleft}. If $n_f = n'_f+1$, by cases on $n_g$
(\nm{onG}). If $n_g = 0$, apply Lemma~\ref{lem:zeroleft} to $g,f$ and
exchange. If $n_g = n'_g+1$, apply Lemma~\ref{lem:tops}.
\end{proof}

\section{The degree (\file{Degree})}

Let $A$ be a gcd domain (module \nm{Over}). To apply
Theorem~\ref{thm:elim} we need a bound on the degree of an element of
$A[X]$. A list gives a bound, its length, but different lists for the
same polynomial give different lengths. So we use the \emph{exact} bound,
which is unique.

\begin{definition}
\begin{itemize}
\item (\nm{TopF}) For $n=0$ the condition is empty (\nm{Top}); for $n+1$
  it is $e_n\neq0$.
\item (\nm{DegF}) A \emph{degree} of $e$ is the datum of $n$ such that
  $e$ is bounded by $n$ and $\nm{TopF}\ n\ e$.
\end{itemize}
\end{definition}

\noindent So the degree in this sense is the usual degree plus one, and
$0$ for the zero polynomial.

\begin{lemma}\label{lem:degprop}
\begin{enumerate}
\item (\nm{top-eq}, \nm{top-prop}) $\nm{TopF}\ n\ e$ is a proposition.
\item (\nm{bnd-prop}) ``$e$ is bounded by $n$'' is a proposition.
\item (\nm{top-vs-bound}) If $e$ is bounded by $m$, then
  $\nm{TopF}\ n\ e$ fails for $n = m+d+1$.
\item (\nm{unique-step}, \nm{unique-cmp}) Two degrees of $e$ are equal.
\item (\nm{deg-prop}) The type of degrees of $e$ is a proposition.
\end{enumerate}
\end{lemma}
\begin{proof}
(1) \nm{Top} has one element, and a negation is a proposition.
(2) A product of equalities in a set.
(3) $n = m+d+1$ is not $0$; and $e_{m+d}=0$ contradicts $e_{m+d}\neq0$.
(4) Let $n$ and $m$ be degrees. Compare them (\nm{unique-cmp}): say
$n = m + d$ (\nm{unique-step}). If $d=0$, then $n=m$. If $d$ is a
successor, (3) gives a contradiction. The other case is symmetric.
(5) By (4), the two first components are equal; the second components are
equal, by (1) and (2).
\end{proof}

\begin{construction}[degree of a list, \nm{deg-list}]
Every list has a degree.
\end{construction}
\begin{proof}
By induction on the list. The empty list has degree $0$. For $c+Xp$, let
$n$ be the degree of $p$ (\nm{cons-deg}). If $n = m+1$, then $c+Xp$ has
degree $m+2$. If $n = 0$, decide whether $c=0$ (\nm{byC}). If $c=0$, then
$c+Xp$ has degree $0$ (\nm{zero-cons}). Otherwise it has degree $1$.
\end{proof}

\begin{definition}[\nm{deg}, \nm{dg}, \nm{dg-bnd}]
Every $x\in A[X]$ has a degree (\nm{deg}): since this is a proposition
(Lemma~\ref{lem:degprop}(5)), it is enough to give it for lists
(\nm{induction}). We write $\deg x$ for the number (\nm{dg}). $x$ is
bounded by $\deg x$ (\nm{dg-bnd}). We let $\cont(x) = \cont_{\deg x}(x)$,
for which Lemma~\ref{lem:bounded} applies (module \nm{Cx}, with \nm{K}).
\end{definition}

\begin{lemma}[\nm{zero-of-coefs}]
If all coefficients of $x$ are $0$, then $x = 0$.
\end{lemma}
\begin{proof}
By \nm{coef-ext}.
\end{proof}

\section{$A[X]$ is a gcd domain (\file{PolyGcd})}

Let $A$ be a gcd domain (module \nm{Over}). For $x\in A[X]$, $\cont(x)$
(\nm{content}) is the content of the coefficients of $x$, bounded by
$\deg x$.

\subsection{The content}

\begin{lemma}\label{lem:contzero}
\begin{enumerate}
\item (\nm{zero-content}) If $\cont(x) = 0$ then $x=0$.
\item (\nm{content-zero}) If all coefficients of $x$ are $0$ then
  $\cont(x)=0$.
\item (\nm{zero-coef}) If $x=0$ then all coefficients of $x$ are $0$.
\item (\nm{zero-mul-coef}) If $h=0$ then all coefficients of $hr$ are
  $0$.
\item (\nm{nz-div}) If $h\mid f$ and $\cont(f)\neq0$ then
  $\cont(h)\neq0$.
\item (\nm{shuffle}) $(cE)q = E(cq)$.
\end{enumerate}
\end{lemma}
\begin{proof}
(1) By \nm{c-zero} and \nm{zero-of-coefs}.
(2) $0$ divides all coefficients, so $0\mid\cont(x)$.
(3) The coefficients of $0$ are $0$.
(4) $hr = 0\cdot r = 0$.
(5) If $\cont(h) = 0$ then $h=0$ by (1); then $f = hr$ has all
coefficients $0$ by (4), and $\cont(f)=0$ by (2).
(6) Associativity and commutativity.
\end{proof}

\begin{lemma}[primitive part, \nm{PP}]\label{lem:pp}
Let $\cont(x)\neq0$, and let $c = \cont(x)$ (\nm{c}; Lemma~\ref{lem:bounded}
applies as \nm{K}, and \nm{S} is the division by $c$). There is $x^*$
(\nm{xs}) with $x = cx^*$ (\nm{eq}), and $x^*$ is primitive (\nm{prim}).
\end{lemma}
\begin{proof}
$c$ divides every coefficient of $x$ (\nm{c-div}), so $x = cx^*$
(\nm{divide}). Let $d$ be a common divisor of $x^*$. Then
$cd\mid cx^*_i = x_i$ for every $i$, so $cd\mid c = c\cdot1$
(\nm{c-univ}). Cancel $c$: $d\mid1$.
\end{proof}

\subsection{The gcd of two non-zero polynomials (module \nm{Main})}

Let $f,g$ have non-zero contents.
\begin{enumerate}
\item (\nm{el}) Let $(D,\lambda)$ be an elimination datum for $f,g$
  (Theorem~\ref{thm:elim}, with the bounds $\deg f$ and $\deg g$); write
  \nm{lam} for $\lambda$.
\item (\nm{nD}) $\cont(D)\neq0$. Otherwise $D = 0$, so
  $\lambda f = Dr = 0$, so $\lambda f_i = 0$ and $f_i = 0$ for every $i$
  ($\lambda\neq0$, \nm{cancel-zero}). Then $\cont(f) = 0$, a
  contradiction.
\item (\nm{PD}, \nm{E}, \nm{cD}) Write $D = \cont(D)E$ with $E$ primitive
  (Lemma~\ref{lem:pp}).
\item (\nm{ef}) $\lambda f = E(\cont(D)r)$, where $\lambda f = Dr$:
  indeed $Dr = (\cont(D)E)r = E(\cont(D)r)$ (\nm{shuffle}).
  (\nm{eg}) Likewise $\lambda g = E(\cont(D)s)$.
\item (\nm{Ef}, \nm{Eg}) $E\mid f$ and $E\mid g$, by
  Lemma~\ref{lem:divprim}.
\item (\nm{gam}) Let $\gamma = \gcdd(\cont(f),\cont(g))$.
  (\nm{ngam}) $\gamma\neq0$, since $\gamma\mid\cont(f)\neq0$.
\item (\nm{Gp}) Let $G = \gamma E$.
\end{enumerate}

\begin{lemma}[\nm{Part}]\label{lem:mainpart}
Let $x\in A[X]$ with $E\mid x$, say $x = Eq$, and $\gamma\mid\cont(x)$.
Then $G\mid x$.
\end{lemma}
\begin{proof}
\begin{itemize}
\item (\nm{dgi}) $\gamma$ divides every coefficient of $Eq$: it divides
  $\cont(x)$, which divides every coefficient of $x = Eq$.
\item (\nm{s}) By Gauss's lemma (Theorem~\ref{thm:gauss}, $E$
  primitive, $\gamma\neq0$), $\gamma$ divides every coefficient of $q$, so
  $q = \gamma s$.
\item (\nm{result}) $x = Eq = E(\gamma s) = (\gamma E)s = Gs$.\qedhere
\end{itemize}
\end{proof}

\begin{lemma}[\nm{Common}]\label{lem:maincommon}
If $h\mid f$ and $h\mid g$ then $h\mid G$.
\end{lemma}
\begin{proof}
\begin{itemize}
\item (\nm{hD}) $h\mid D$, by the elimination datum.
\item (\nm{nh}) $\cont(h)\neq0$, by Lemma~\ref{lem:contzero}(5).
\item (\nm{PH}, \nm{ch}) Write $h = \cont(h)h^*$ with $h^*$ primitive
  (Lemma~\ref{lem:pp}).
\item (\nm{ch-div}) If $h\mid x$ then $\cont(h)$ divides every
  coefficient of $x$: $x = hr = \cont(h)(h^*r)$.
\item (\nm{ch-gam}) $\cont(h)\mid\gamma$: $\cont(h)$ divides $\cont(f)$
  and $\cont(g)$, by (\nm{ch-div}) and \nm{c-univ}.
\item (\nm{eD}) $\cont(D)E = D = ht = (\cont(h)h^*)t = h^*(\cont(h)t)$.
\item (\nm{hsE}) $h^*\mid E$, by Lemma~\ref{lem:divprim} with
  $\lambda = \cont(D)\neq0$.
\item (\nm{cC}) $\cont(h)\mid\gamma$ also holds for the constants in
  $A[X]$.
\item (\nm{result}) $h = \cont(h)h^*\mid\gamma E = G$
  (\nm{div-mul2}).\qedhere
\end{itemize}
\end{proof}

\noindent (\nm{gcdOf}) Hence $G$ is a gcd of $f$ and $g$:
$G\mid f$ and $G\mid g$ by Lemma~\ref{lem:mainpart} (with
$\gamma\mid\cont(f)$, resp.\ $\gamma\mid\cont(g)$), and every common
divisor divides $G$ by Lemma~\ref{lem:maincommon}.

\begin{theorem}[\nm{gcdP}]\label{thm:gcdP}
Any two $f,g\in A[X]$ have a gcd.
\end{theorem}
\begin{proof}
Decide whether $\cont(f) = 0$ (\nm{byF}). If so, $f = 0$, and $g$ is a
gcd. Otherwise decide whether $\cont(g) = 0$ (\nm{byG}). If so, $g = 0$,
and $f$ is a gcd. Otherwise apply \nm{Main}.
\end{proof}

\subsection{Discrete, non-trivial, integral}

\begin{lemma}
\begin{enumerate}
\item (\nm{decP}, by cases \nm{byDec}) $x = 0$ or $x\neq0$.
\item (\nm{nontrivialP}) $1\neq0$ in $A[X]$.
\end{enumerate}
\end{lemma}
\begin{proof}
(1) Decide whether $\cont(x)=0$. If so, $x = 0$
(Lemma~\ref{lem:contzero}(1)). If not, $x\neq0$, since $x = 0$ would give
$\cont(x) = 0$ (Lemma~\ref{lem:contzero}(3),(2)).
(2) The constant coefficient of $1$ is $1$, and that of $0$ is $0$.
\end{proof}

\begin{lemma}[integrality, \nm{integralP}]
If $x\neq0$ and $y\neq0$ then $xy\neq0$.
\end{lemma}
\begin{proof}
Let \nm{Integral}~$x\,y$ be the statement. It is a proposition, so it is
enough to prove it for classes of lists $p$, $q$ (\nm{forall2},
\nm{integral-lists}).
\begin{itemize}
\item (\nm{list-content-nz}) A non-zero list has non-zero content: if
  $\cont(p) = 0$ then all coefficients of $p$ are $0$, so its class is
  $0$.
\item Suppose $pq = 0$. Then $0$ divides every coefficient of $pq$, so
  $0\mid\cont(p)\cont(q)$ by Theorem~\ref{thm:C}, that is,
  $\cont(p)\cont(q) = 0$. This contradicts integrality of $A$, since
  $\cont(p)\neq0$ and $\cont(q)\neq0$.\qedhere
\end{itemize}
\end{proof}

\begin{theorem}[\nm{domain}]
If $A$ is a gcd domain, then $A[X]$ is a gcd domain.
\end{theorem}
\begin{proof}
$A[X]$ is a set (\nm{set}), and it is discrete (\nm{decP}), integral
(\nm{integralP}), has gcds (\nm{gcdP}) and $1\neq0$ (\nm{nontrivialP}).
\end{proof}

\section*{What a usual text leaves implicit}

A comparison of this text with MRR shows the steps that the formal proof
makes explicit and an informal proof usually omits.
\begin{enumerate}
\item \emph{Gcds and quotients are data.} A divisibility $a\mid b$ is a
  quotient with an equation. A gcd is a chosen element with its
  properties. This is why $\cont(fg)$ and $\cont(f)\cont(g)$ are compared
  by divisibility both ways, and never by equality.
\item \emph{Reasoning on a representative.} $A[X]$ is a quotient of lists.
  Each time a statement is proved ``on the lists'', the text checks that
  it is a proposition (\nm{quot-prop}, \nm{stmt-prop}, \nm{deg-prop},
  \nm{integralP}). This is why Theorem~\ref{thm:gauss} assumes $d\neq0$.
  It is also why the degree must be the exact one, and not the length of
  some list.
\item \emph{The induction hypothesis} of Gauss's lemma is the general
  statement (C), not the primitive case (Section~\ref{sec:gauss},
  point~3).
\item \emph{Cases on decisions.} Every ``if $x=0$ \dots, otherwise
  \dots'' uses the decision \nm{dec} of the gcd domain. It is written as
  its own small declaration (\nm{lemma-M-byU}, \nm{step-byP},
  \nm{byF}, \dots).
\item \emph{Rewriting.} Steps such as ``$d\mid b$ and $b=b'$, so
  $d\mid b'$'' (\nm{div-eq-right}) and index arithmetic ($n'_f+(n'_g+1)=N$,
  \nm{eF}) appear as explicit steps.
\end{enumerate}

\begin{thebibliography}{9}
\bibitem{jutting} L.~S.~van Benthem Jutting, \emph{The development of a
  text in AUT-QE}. Preprint, Technological University Eindhoven; in
  P.~Braffort (ed.), Proceedings of the Symposium on APL, Paris,
  December 1973.
\bibitem{mrr} R.~Mines, F.~Richman and W.~Ruitenburg, \emph{A Course in
  Constructive Algebra}. Springer, 1988.
\end{thebibliography}

\end{document}
